\chapter{Introducción}
\label{cap:introduccion}

\section{Motivación}

El lenguaje es una de las capacidades más complejas y distintivas de la especie humana. En el mundo existen miles de lenguas vivas \citep{eberhard2020ethnologue}, pero la gran mayoría dispone de una presencia digital muy reducida, en muchos casos apenas existen corpus textuales y, por consiguiente, herramientas de procesamiento automático. Esta escasez de información hace que los sistemas de inteligencia artificial que dependen de grandes volúmenes de datos de entrenamiento, tengan muchas dificultades para trabajar con ellas, al menos de forma fiable.

Sin embargo, las personas somos capaces de aprender a leer, escribir y traducir en un idioma completamente desconocido a partir de un número reducido de ejemplos, gracias a nuestra \textit{competencia metalingüística}. En otras palabras, podemos identificar el vocabulario, reconocer patrones morfológicos y formular reglas gramaticales únicamente comparando palabras y frases. Esta capacidad más característica y propia de los seres humanos plantea una pregunta fundamental para la inteligencia artificial: ¿pueden los modelos de lenguaje de gran escala hacer lo mismo?

\textit{La Olimpiada Internacional de Lingüística (IOL)} propone anualmente problemas del tipo \textit{piedra Rosetta}: se presenta al participante un conjunto de frases en una lengua completamente desconocida junto con sus traducciones, y se le pide que deduzca las reglas del idioma y las aplique para traducir nuevas frases \citep{derzhanski2010linguistics,bozhanov2013rosetta,iol2026problems}. Este tipo de problemas son autocontenidos, es decir, toda la información necesaria se encuentra en el enunciado, y requieren razonamiento explícito, no memorización. Por la forma en la que se presenta la información, los ejercicios de las olimpiadas lingüísticas presentan un banco de pruebas ideal para medir si un sistema de inteligencia artificial posee razonamiento metalingüístico genuino, o al menos si es capaz de reproducir este tipo de razonamiento \citep{sahin2020puzzling,sanchez2024linguini}.

En este trabajo vamos a partir de la hipótesis de que es posible diseñar estrategias de prompting que permitan a los modelos de lenguaje resolver estos problemas de forma progresiva, construyendo explícitamente un diccionario y un conjunto de reglas gramaticales a partir de los ejemplos proporcionados, sin recibir ninguna pista sobre la lengua en cuestión. En el caso de que la hipótesis se confirme, esto abriría la puerta a utilizar modelos de lenguaje para procesar, documentar y analizar miles de lenguas del mundo para las que no existe ni existirá suficiente cantidad de datos de entrenamiento. En caso contrario, estudiaremos y evidenciaremos hasta qué punto los modelos actuales pueden reproducir un proceso de análisis similar al que sigue una persona al enfrentarse a una lengua que no conoce.

Cabe mencionar que la motivación del proyecto también está relacionada con mi interés personal por los idiomas y por la forma en que distintas lenguas expresan una misma realidad mediante formas tan diferentes.

\section{Contexto y estado del problema}

Durante los últimos años, los modelos de lenguaje de gran escala han demostrado capacidades notables en tareas como la generación de texto, la traducción automática o la respuesta a preguntas \citep{brown2020language}. Una de sus características más relevantes es el aprendizaje en contexto, que es la posibilidad de realizar una tarea a partir de ejemplos incluidos directamente en el prompt, es decir, en la forma en la que se le presenta la información al modelo mediante la petición, sin necesidad de modificar sus parámetros. Aquí vamos a plantear una pregunta especialmente relevante para las lenguas de baja densidad de recursos: ¿puede un modelo deducir patrones lingüísticos a partir de unos pocos ejemplos, incluso cuando no dispone de información suficiente para apoyarse únicamente en los datos aprendidos durante su entrenamiento?

En los últimos años se han propuesto varios conjuntos de evaluación basados en problemas de olimpiadas lingüísticas, de los que hablaremos con más detalle en el Capítulo~\ref{cap:marcoteorico}. Estos datasets reúnen problemas de distintas lenguas y permiten evaluar la capacidad de los modelos para resolver tareas como la identificación de patrones morfológicos y la traducción de textos \citep{sahin2020puzzling,chi2024modeling,sanchez2024linguini,bean2024lingoly}. Los resultados publicados hasta el momento muestran que estos problemas siguen siendo difíciles para los modelos actuales. Aunque algunos sistemas obtienen respuestas parcialmente correctas, su rendimiento se mantiene por debajo del nivel humano en muchos tipos de ejercicios, especialmente cuando las lenguas presentan una morfología compleja o estructuras alejadas de las más frecuentes en los datos de entrenamiento.

Distintos trabajos han estudiado formas de mejorar estos resultados mediante estrategias de prompting más elaboradas. En lugar de solicitar directamente una traducción final, estas estrategias guían al modelo para que identifique vocabulario, deduzca reglas y revise sus hipótesis antes de responder \citep{wei2022chain,zhu2025evaluating,madaan2023selfrefine}. Precisamente de esto trata este trabajo de fin de grado. El objetivo es diseñar, implementar y comparar varias estrategias de prompting para la resolución de problemas de olimpiadas lingüísticas. En particular, se estudiarán distintas maneras de presentación de los ejemplos al modelo, que a su vez puedan permitir construir conocimiento de forma progresiva y corregir hipótesis incorrectas a medida que recibe nueva información.

\section{Estructura del documento}

El presente documento se organiza en los siguientes capítulos:

El \textbf{Capítulo~\ref{cap:marcoteorico}} presenta el estado de la cuestión, revisando los fundamentos teóricos del razonamiento metalingüístico, los principales benchmarks de olimpiadas lingüísticas para LLMs y las estrategias de prompting más relevantes de la literatura.

El \textbf{Capítulo~\ref{cap:objetivos}} define el objetivo principal del trabajo y los objetivos específicos que se derivan de él, así como las restricciones y condicionantes del proyecto.

El \textbf{Capítulo~\ref{cap:metodologia}} describe la metodología del trabajo: los puzzles utilizados, los modelos evaluados, el control de dataset previo, las estrategias de prompting diseñadas, el sistema experimental implementado y las métricas de evaluación empleadas.

El \textbf{Capítulo~\ref{cap:desarrollo}} detalla la implementación del sistema experimental y las decisiones adoptadas durante su desarrollo.

El \textbf{Capítulo~\ref{cap:resultados}} presenta y analiza los resultados experimentales obtenidos, incluyendo la comparación con los valores de referencia del benchmark \textsc{Linguini} \citep{sanchez2024linguini} y el análisis cualitativo comparando el razonamiento del modelo con el humano. Se incluyen pruebas adicionales y complementarias que ayuden a evaluar el comportamiento del modelo para diferentes tipos de problemas.

El \textbf{Capítulo~\ref{cap:conclusiones}} recoge las conclusiones principales del trabajo, destacando los resultados más relevantes, las aportaciones realizadas y las limitaciones observadas durante la evaluación.
 
Finalmente, el \textbf{Capítulo~\ref{cap:trabajos_futuros}} plantea las líneas de trabajo futuro que se abren a partir de los resultados obtenidos.